\documentclass[11pt]{article}

\usepackage[utf8]{inputenc}
\usepackage[T1]{fontenc}
\usepackage{lmodern}
\usepackage{geometry}
\usepackage{amsmath,amssymb,amsfonts,amsthm,mathtools}
\usepackage{bm}
\usepackage{enumitem}
\usepackage{microtype}
\usepackage{hyperref}
\usepackage{titlesec}
\usepackage{booktabs}
\usepackage{array}
\usepackage{setspace}
\usepackage{mathrsfs}
\usepackage{physics}

\geometry{margin=1in}
\hypersetup{colorlinks=true, linkcolor=black, urlcolor=blue, citecolor=black}
\setlist[enumerate]{topsep=0.4em,itemsep=0.35em}
\setlist[itemize]{topsep=0.3em,itemsep=0.25em}
\titleformat{\section}{\large\bfseries}{\thesection.}{0.5em}{}
\titleformat{\subsection}{\normalsize\bfseries}{\thesubsection.}{0.5em}{}
\setstretch{1.06}

\newcommand{\Aether}{\AE{}ther}
\newcommand{\Aflow}{\AE{}ther-flow}
\newcommand{\TheoryName}{\Aflow{} Interpretation of Relativity}
\newcommand{\TheoryProgram}{\Aether{} / \Aflow{} framework}

\newcommand{\CompletedMetric}{g^{\sharp}}
\newcommand{\MatterSpace}{\Psi}

\newcommand{\Car}{\mathrm{car}}
\newcommand{\Cur}{\mathrm{cur}}
\newcommand{\Hom}{\mathrm{hom}}

\newcommand{\ForSpace}[1]{\mathcal I_{#1}^{\mathrm{for}}}
\newcommand{\PrimShell}{\mathcal S_{\mathrm{prim}}}
\newcommand{\Reduction}[1]{\mathcal R_{#1}}
\newcommand{\CurrentCarrier}[1]{\mathfrak C_{#1}^{\mathrm{cur}}}
\newcommand{\EqCarrier}[1]{\mathfrak C_{#1}^{\mathrm{eq}}}
\newcommand{\MetricOut}{\mathscr G_{\mathrm{out}}}
\newcommand{\MatterOut}{\mathscr T_{\mathrm{out}}}

\newcommand{\SubTarget}[1]{E_{#1}^{\mathrm{sub}}}
\newcommand{\EqnTarget}[1]{Q_{#1}^{\mathrm{eqn}}}
\newcommand{\HiddenSpace}[1]{\mathcal H_{#1}^{\mathrm{hid}}}
\newcommand{\HiddenBody}[1]{D_{#1}^{\mathrm{hid}}}

\newcommand{\SeedState}[1]{\mathcal S_{#1}^{\mathrm{seed}}}
\newcommand{\SeedBody}[1]{\mathcal B_{#1}^{\mathrm{seed}}}
\newcommand{\SeedShell}[1]{\mathcal Z_{#1}^{\mathrm{seed}}}
\newcommand{\SeedSupportShell}[1]{\mathcal U_{#1}^{\mathrm{seed,sup}}}
\newcommand{\SeedZeroCore}[1]{\mathcal Z_{#1}^{0,\mathrm{seed}}}
\newcommand{\SeedSplit}[1]{\Phi_{#1}^{\mathrm{seed}}}
\newcommand{\SeedCharge}[1]{\mathcal Q_{#1}^{\mathrm{seed}}}
\newcommand{\SeedEmbed}[1]{\zeta_{#1}^{\mathrm{seed}}}
\newcommand{\SeedKernel}[1]{\widetilde K_{#1}^{0,\mathrm{seed}}}
\newcommand{\SeedCoercive}[1]{P_{#1}^{\mathrm{seed,co}}}

\newcommand{\GermTransportCore}{\mathfrak G^{\mathrm{sub,germ,stc}}}
\newcommand{\GermNucleus}{\mathfrak G^{\mathrm{sub,germ,zstn}}}
\newcommand{\GermState}[1]{\mathcal G_{#1}^{\mathrm{germ}}}
\newcommand{\GermBody}[1]{\mathcal B_{#1}^{\mathrm{germ}}}
\newcommand{\GermProj}[1]{\Pi_{#1}^{\mathrm{germ}}}
\newcommand{\GermSeedSec}[1]{\sigma_{#1}^{\mathrm{germ\leftarrow seed}}}
\newcommand{\GermSection}[1]{\Gamma_{#1}^{\mathrm{germ}}}
\newcommand{\GermShell}[1]{\mathcal Z_{#1}^{\mathrm{germ}}}
\newcommand{\GermSupportShell}[1]{\mathcal U_{#1}^{\mathrm{germ,sup}}}
\newcommand{\GermZeroCore}[1]{\mathcal Z_{#1}^{0,\mathrm{germ}}}
\newcommand{\GermSplit}[1]{\Phi_{#1}^{\mathrm{germ}}}
\newcommand{\GermCharge}[1]{\mathcal Q_{#1}^{\mathrm{germ}}}
\newcommand{\GermEmbed}[1]{\zeta_{#1}^{\mathrm{germ}}}
\newcommand{\GermKernel}[1]{\widetilde K_{#1}^{0,\mathrm{germ}}}
\newcommand{\GermCoercive}[1]{P_{#1}^{\mathrm{germ,co}}}

\newtheorem{definition}{Definition}
\newtheorem{proposition}{Proposition}
\newtheorem{remark}{Remark}
\newtheorem{corollary}{Corollary}
\newtheorem{theorem}{Theorem}

\title{\textbf{The \TheoryName{}}\\[0.4em]
\large Primitive-Reservoir Observer-Reduction\\
\large Observer-Relevant Germ Seed-Transport Core\\
\large Collapse to a Minimal Germ\\
\large Zero-Shell Transport Nucleus}
\author{}
\date{}

\begin{document}

\maketitle

\begin{abstract}
The current live benchmark-facing gain on the active primitive-reservoir observer-reduction line is the theorem collapsing the full seed-generating substrate germ package to the smaller observer-relevant germ seed-transport core \(\GermTransportCore\). Under the recorded routing safeguard, the next honest move below that frontier had to do one of two things: derive or uniquely force \(\GermTransportCore\) from still more primitive explicit substrate structure, or prove a genuinely smaller theorem-level collapse strictly inside \(\GermTransportCore\). The present note takes that admissible fallback route.

For each sector it isolates a smaller \emph{minimal germ zero-shell transport nucleus} \(\GermNucleus\). The nucleus keeps only the germ state space and compact germ body, the projection to the fixed seed body together with its seed section, the germ restriction section, the exact germ shell, and benchmark faithfulness. It omits, as independent benchmark-facing data, the germ support shell, the germ zero core, the germ hidden split, the hidden charge, the exact germ zero-response realization, and the fixed-data solvability / coercive package.

The main theorem shows that those omitted constituents are canonically induced once the nucleus is fixed against the recorded downstream seed package. The support shell and zero core are inverse images of the fixed seed support shell and seed zero core under the exact-shell transport diffeomorphism; the hidden split is the lifted fixed seed split; the hidden charge is the fixed seed charge itself; and the exact germ zero-response realization together with the solvability / coercive package are transported uniquely across the canonical zero-core diffeomorphism. Hence the current germ seed-transport core carries no additional benchmark-facing freedom beyond the smaller nucleus. The benchmark boundary remains unchanged throughout: the one operative metric is still exactly \(\CompletedMetric\), the ordinary matter route is unchanged, there is no second operative metric, and the low-energy matter law is not enlarged. The positive result is therefore narrow but real: the live stopping point shrinks further from the observer-relevant germ seed-transport core to a smaller minimal germ zero-shell transport nucleus.
\end{abstract}

\section{Introduction}

The active derivational program of \TheoryName{} is no longer at a stage where another same-output deeper-origin relay below the current germ-core frontier can count merely by going one layer deeper. The current germ-core collapse theorem already spent the honest first internal compression move available there: once the fixed downstream seed package is held in hand, the lower minimizer ladder beneath seed scope is no longer an independent benchmark-facing burden. What remained live after that theorem was the observer-relevant germ seed-transport core \(\GermTransportCore\).

The next honest options were therefore narrowed again. One could either derive or uniquely force \(\GermTransportCore\) from still more primitive explicit substrate structure in a way that genuinely changes benchmark-facing status, or ask whether \(\GermTransportCore\) itself still contains benchmark-neutral constituents and therefore collapses further. The present note takes the second route. That choice is disciplined. Until one checks whether the current core is already minimal, a merely deeper-origin manuscript that reproduces the same core risks becoming another same-output relay.

The sharper theorem proved here is that several current-core constituents are already forced by exact-shell transport to the fixed seed package. Once one fixes the germ state space and body, the projection to the fixed seed body together with its seed section, the germ restriction section, the exact germ shell, and benchmark faithfulness, the remaining current-core data no longer carry independent benchmark-facing freedom. The germ support shell and germ zero core are the inverse images of the fixed seed support shell and seed zero core; the hidden split is the lifted fixed seed split; the hidden charge is the fixed seed charge; the exact zero-response realization is the lifted fixed seed realization; and the fixed-data solvability / coercive package is transported uniquely from the fixed seed package across the canonical zero-core diffeomorphism.

\paragraph{Framework and claim boundary.}
\TheoryProgram{} denotes the broader \Aether{} / \Aflow{} framework built on the claims that the \Aether{} is the underlying four-dimensional substrate of reality, the \Aflow{} is its intrinsic ordered motion, observed three-dimensional space is the local experiential slice of that deeper substrate, S-time is the experienced order of change arising from matter, light, and the \Aflow{}, and observed expansion is the three-dimensional appearance of deeper four-dimensional ordered motion. Within that framework, \TheoryName{} denotes the exact-closure theory in which the effective gravitational dynamics are adopted to be exactly Einsteinian with universal matter coupling. In the active sequence, \emph{adoption} means use of that established relativistic dynamics without claiming substrate derivation, while \emph{derivation} is reserved for a first-principles recovery from explicit substrate variables.

The present note stays strictly inside that boundary. It does not alter the benchmark exact-closure package. It does not introduce a second operative metric or a new low-energy matter law. It only proves that the current observer-relevant germ seed-transport core collapses further to a smaller minimal germ zero-shell transport nucleus.

\section{Fixed Inputs}

\begin{definition}[Recorded primitive continuation data]
\label{def:fixed}
Fix the following continuation data from the active primitive-reservoir observer-reduction line.
\begin{enumerate}
    \item For each sector label \(\sigma\in\{\Car,\Cur,\Hom\}\), a primitive forcing shell
    \[
    \ForSpace{\sigma},
    \]
    a visible reduction map
    \[
    \Reduction{\sigma}:\ForSpace{\sigma}\to X_{\sigma},
    \]
    and shell-level current and equilibrium carriers
    \[
    \CurrentCarrier{\sigma}:\ForSpace{\sigma}\to Z_{\sigma}^{\mathrm{cur}},
    \qquad
    \EqCarrier{\sigma}:\ForSpace{\sigma}\to Z_{\sigma}^{\mathrm{eq}}.
    \]
    \item The product shell
    \[
    \PrimShell
    \equiv
    \ForSpace{\Car}\times \ForSpace{\Cur}\times \ForSpace{\Hom}.
    \]
    \item The recorded metric and ordinary-matter outputs
    \[
    \MetricOut:\PrimShell\to\{\text{observer metrics}\},
    \qquad
    \MatterOut:\PrimShell\times\MatterSpace\to\{\text{observer stress tensors}\},
    \]
    satisfying
    \begin{equation}
    \MetricOut(\mathbf w)=\CompletedMetric,
    \qquad
    \MatterOut(\mathbf w,\psi)
    =
    T_{\mu\nu}^{\mathrm{matter}}[\CompletedMetric,\psi]
    \label{eq:fixedoutputs}
    \end{equation}
    for every \(\mathbf w\in\PrimShell\) and every matter configuration \(\psi\in\MatterSpace\).
\end{enumerate}
\end{definition}

\begin{definition}[Recorded downstream seed package]
\label{def:seed}
Fix \(\sigma\in\{\Car,\Cur,\Hom\}\). The active line already fixes the downstream seed package at current benchmark-facing scope:
\begin{enumerate}
    \item a finite-dimensional seed state space \(\SeedState{\sigma}\), a compact convex seed body \(\SeedBody{\sigma}\subset\SeedState{\sigma}\), the exact seed shell \(\SeedShell{\sigma}\subset\SeedBody{\sigma}\), and the exact seed support shell \(\SeedSupportShell{\sigma}\subset\SeedShell{\sigma}\);
    \item a seed zero core \(\SeedZeroCore{\sigma}\subset\SeedSupportShell{\sigma}\), the same hidden fiber space \(\HiddenSpace{\sigma}\) with compact hidden body \(\HiddenBody{\sigma}\subset\HiddenSpace{\sigma}\), and the fixed seed split
    \[
    \SeedSplit{\sigma}:
    \SeedZeroCore{\sigma}\times\HiddenBody{\sigma}
    \xrightarrow{\cong}
    \SeedSupportShell{\sigma};
    \]
    \item the seed hidden charge
    \[
    \SeedCharge{\sigma}:\HiddenSpace{\sigma}\to\EqnTarget{\sigma},
    \]
    a smooth observer-compatible exact seed zero-response realization
    \[
    \SeedEmbed{\sigma}:\ForSpace{\sigma}\hookrightarrow\SeedZeroCore{\sigma},
    \]
    and the fixed-data solvability / coercive package recorded through \(\SeedKernel{\sigma}\) and \(\SeedCoercive{\sigma}\);
    \item benchmark faithfulness, meaning preservation of the benchmark outputs \eqref{eq:fixedoutputs}, no second operative metric, no enlarged low-energy matter law, and no premature promotion from adoption to completed derivation.
\end{enumerate}
\end{definition}

\begin{remark}
Definitions \ref{def:fixed} and \ref{def:seed} are fixed inputs. The primitive continuation data, the one operative metric, the ordinary matter route, and the recorded downstream seed package are not reopened here. The present note asks only whether the current observer-relevant germ seed-transport core is already minimal once those downstream data are held fixed.
\end{remark}

\section{Current Observer-Relevant Germ Seed-Transport Core}

\begin{definition}[Observer-relevant germ seed-transport core]
\label{def:core}
Fix \(\sigma\in\{\Car,\Cur,\Hom\}\). The \emph{observer-relevant germ seed-transport core} \(\GermTransportCore\) for sector \(\sigma\) consists of:
\begin{enumerate}
    \item the germ state space \(\GermState{\sigma}\), compact convex germ body \(\GermBody{\sigma}\subset\GermState{\sigma}\), projection
    \[
    \GermProj{\sigma}:\GermState{\sigma}\to\SeedState{\sigma},
    \]
    and seed section
    \[
    \GermSeedSec{\sigma}:\SeedBody{\sigma}\to\GermBody{\sigma}
    \]
    satisfying
    \[
    \GermProj{\sigma}\bigl(\GermBody{\sigma}\bigr)=\SeedBody{\sigma},
    \qquad
    \GermProj{\sigma}\circ\GermSeedSec{\sigma}
    =
    \mathrm{id}_{\SeedBody{\sigma}};
    \]
    \item the restriction section
    \[
    \GermSection{\sigma}:\GermState{\sigma}\to\SubTarget{\sigma},
    \]
    and the exact germ shell
    \[
    \GermShell{\sigma}
    \equiv
    \bigl(\GermSection{\sigma}\bigr)^{-1}(0)\cap\GermBody{\sigma},
    \]
    with projected exact-shell transport
    \[
    \GermProj{\sigma}\big|_{\GermShell{\sigma}}
    :
    \GermShell{\sigma}
    \xrightarrow{\cong}
    \SeedShell{\sigma};
    \]
    \item the germ support shell \(\GermSupportShell{\sigma}\subset\GermShell{\sigma}\), the germ zero core \(\GermZeroCore{\sigma}\subset\GermSupportShell{\sigma}\), and the germ split
    \[
    \GermSplit{\sigma}:
    \GermZeroCore{\sigma}\times\HiddenBody{\sigma}
    \xrightarrow{\cong}
    \GermSupportShell{\sigma}
    \]
    satisfying
    \[
    \GermProj{\sigma}\big|_{\GermSupportShell{\sigma}}
    :
    \GermSupportShell{\sigma}
    \xrightarrow{\cong}
    \SeedSupportShell{\sigma},
    \qquad
    \GermProj{\sigma}\big|_{\GermZeroCore{\sigma}}
    :
    \GermZeroCore{\sigma}
    \xrightarrow{\cong}
    \SeedZeroCore{\sigma},
    \]
    and
    \[
    \SeedSplit{\sigma}(z,\eta)
    =
    \GermProj{\sigma}\Bigl(
    \GermSplit{\sigma}\bigl(
    (\GermProj{\sigma}\big|_{\GermZeroCore{\sigma}})^{-1}(z),\eta
    \bigr)\Bigr)
    \]
    for every \(z\in\SeedZeroCore{\sigma}\) and every \(\eta\in\HiddenBody{\sigma}\);
    \item the hidden charge \(\GermCharge{\sigma}\), the exact germ zero-response realization
    \[
    \GermEmbed{\sigma}:\ForSpace{\sigma}\hookrightarrow\GermZeroCore{\sigma},
    \]
    and the fixed-data solvability / coercive package \(\GermKernel{\sigma},\GermCoercive{\sigma}\), all projecting through \(\GermProj{\sigma}\big|_{\GermZeroCore{\sigma}}\) to the fixed seed hidden charge \(\SeedCharge{\sigma}\), exact seed zero-response realization \(\SeedEmbed{\sigma}\), and fixed-data solvability / coercive package \(\SeedKernel{\sigma},\SeedCoercive{\sigma}\);
    \item benchmark faithfulness in the sense of Definition \ref{def:seed}(4).
\end{enumerate}
\end{definition}

\begin{remark}
Definition \ref{def:core} is the current benchmark-facing frontier before the present theorem is applied. The question now is whether every constituent listed there is still an independent live burden, or whether some of those constituents are already forced once the exact-shell transport data in items (1) and (2) are fixed against the recorded downstream seed package.
\end{remark}

\section{Minimal Germ Zero-Shell Transport Nucleus}

\begin{definition}[Minimal germ zero-shell transport nucleus]
\label{def:nucleus}
Fix \(\sigma\in\{\Car,\Cur,\Hom\}\). The \emph{minimal germ zero-shell transport nucleus} \(\GermNucleus\) for sector \(\sigma\) consists of:
\begin{enumerate}
    \item the germ state space \(\GermState{\sigma}\), compact convex germ body \(\GermBody{\sigma}\subset\GermState{\sigma}\), projection
    \[
    \GermProj{\sigma}:\GermState{\sigma}\to\SeedState{\sigma},
    \]
    and seed section
    \[
    \GermSeedSec{\sigma}:\SeedBody{\sigma}\to\GermBody{\sigma}
    \]
    satisfying
    \[
    \GermProj{\sigma}\bigl(\GermBody{\sigma}\bigr)=\SeedBody{\sigma},
    \qquad
    \GermProj{\sigma}\circ\GermSeedSec{\sigma}
    =
    \mathrm{id}_{\SeedBody{\sigma}};
    \]
    \item the restriction section
    \[
    \GermSection{\sigma}:\GermState{\sigma}\to\SubTarget{\sigma},
    \]
    and the exact germ shell
    \[
    \GermShell{\sigma}
    \equiv
    \bigl(\GermSection{\sigma}\bigr)^{-1}(0)\cap\GermBody{\sigma},
    \]
    with
    \[
    \GermProj{\sigma}\big|_{\GermShell{\sigma}}
    :
    \GermShell{\sigma}
    \xrightarrow{\cong}
    \SeedShell{\sigma};
    \]
    \item benchmark faithfulness in the sense of Definition \ref{def:core}(5).
\end{enumerate}
By definition, the current-core constituents
\[
\GermSupportShell{\sigma},
\quad
\GermZeroCore{\sigma},
\quad
\GermSplit{\sigma},
\quad
\GermCharge{\sigma},
\quad
\GermEmbed{\sigma},
\quad
\GermKernel{\sigma},
\quad
\GermCoercive{\sigma}
\]
are \emph{not} taken as independent nucleus data.
\end{definition}

\begin{definition}[Core completion of the nucleus]
\label{def:completion}
A \emph{benchmark-faithful core completion} of \(\GermNucleus\) is an observer-relevant germ seed-transport core in the sense of Definition \ref{def:core} that agrees with the given nucleus and adds only the omitted current-core constituents listed in Definition \ref{def:nucleus}.
\end{definition}

\begin{remark}
Definition \ref{def:nucleus} keeps exactly the data needed to present germ exact-shell transport to the fixed downstream seed package while preserving the benchmark one-metric observer package. The present theorem asks whether the omitted support-shell, zero-core, hidden-sector, realization, and kernel data remain independently benchmark-facing once that smaller nucleus is fixed.
\end{remark}

\section{Canonical Completion from the Nucleus}

\begin{proposition}[Canonical induced support shell]
\label{prop:support}
Fix a minimal germ zero-shell transport nucleus \(\GermNucleus\). Then the omitted germ support shell is determined canonically by
\[
\GermSupportShell{\sigma}
\equiv
\bigl(\GermProj{\sigma}\big|_{\GermShell{\sigma}}\bigr)^{-1}
\bigl(\SeedSupportShell{\sigma}\bigr).
\]
Moreover,
\[
\GermProj{\sigma}\big|_{\GermSupportShell{\sigma}}
:
\GermSupportShell{\sigma}
\xrightarrow{\cong}
\SeedSupportShell{\sigma}
\]
is a diffeomorphism.
\end{proposition}

\begin{proof}
Definition \ref{def:nucleus}(2) already requires
\[
\GermProj{\sigma}\big|_{\GermShell{\sigma}}
:
\GermShell{\sigma}
\xrightarrow{\cong}
\SeedShell{\sigma}.
\]
Because the fixed seed support shell \(\SeedSupportShell{\sigma}\subset\SeedShell{\sigma}\) is part of Definition \ref{def:seed}, its inverse image inside \(\GermShell{\sigma}\) is canonically defined. Since the shell projection is a diffeomorphism, its restriction to that inverse image is again a diffeomorphism onto \(\SeedSupportShell{\sigma}\). No additional choice remains.
\end{proof}

\begin{proposition}[Canonical induced zero core]
\label{prop:zero}
Fix a minimal germ zero-shell transport nucleus \(\GermNucleus\). Then the omitted germ zero core is determined canonically by
\[
\GermZeroCore{\sigma}
\equiv
\bigl(\GermProj{\sigma}\big|_{\GermShell{\sigma}}\bigr)^{-1}
\bigl(\SeedZeroCore{\sigma}\bigr).
\]
Moreover,
\[
\GermZeroCore{\sigma}\subset\GermSupportShell{\sigma}
\]
and
\[
\GermProj{\sigma}\big|_{\GermZeroCore{\sigma}}
:
\GermZeroCore{\sigma}
\xrightarrow{\cong}
\SeedZeroCore{\sigma}
\]
is a diffeomorphism.
\end{proposition}

\begin{proof}
The fixed seed zero core \(\SeedZeroCore{\sigma}\subset\SeedSupportShell{\sigma}\subset\SeedShell{\sigma}\) is part of Definition \ref{def:seed}. Because Definition \ref{def:nucleus}(2) supplies a shell diffeomorphism from \(\GermShell{\sigma}\) to \(\SeedShell{\sigma}\), the inverse image of \(\SeedZeroCore{\sigma}\) is canonically defined inside \(\GermShell{\sigma}\). Since \(\SeedZeroCore{\sigma}\subset\SeedSupportShell{\sigma}\), Proposition \ref{prop:support} shows that this inverse image lies inside \(\GermSupportShell{\sigma}\). The restricted projection is a diffeomorphism because it is the restriction of the shell diffeomorphism.
\end{proof}

\begin{proposition}[Canonical split and hidden charge]
\label{prop:split}
Fix a minimal germ zero-shell transport nucleus \(\GermNucleus\). Then the omitted germ split and hidden charge are determined canonically by
\[
\GermCharge{\sigma}\equiv \SeedCharge{\sigma}
\]
and
\[
\GermSplit{\sigma}(z,\eta)
\equiv
\bigl(\GermProj{\sigma}\big|_{\GermShell{\sigma}}\bigr)^{-1}
\Bigl(
\SeedSplit{\sigma}\bigl(
(\GermProj{\sigma}\big|_{\GermZeroCore{\sigma}})(z),\eta
\bigr)
\Bigr)
\]
for every \(z\in\GermZeroCore{\sigma}\) and every \(\eta\in\HiddenBody{\sigma}\). These data satisfy the projected split identity of Definition \ref{def:core}(3).
\end{proposition}

\begin{proof}
The hidden charge is already fixed on the downstream seed package by Definition \ref{def:seed}(3), so requiring \(\GermCharge{\sigma}=\SeedCharge{\sigma}\) leaves no freedom. For the split, Proposition \ref{prop:zero} supplies the zero-core diffeomorphism
\[
\GermProj{\sigma}\big|_{\GermZeroCore{\sigma}}
:
\GermZeroCore{\sigma}
\xrightarrow{\cong}
\SeedZeroCore{\sigma},
\]
while Definition \ref{def:nucleus}(2) supplies the shell diffeomorphism to \(\SeedShell{\sigma}\). Composing those fixed diffeomorphisms with the fixed seed split \(\SeedSplit{\sigma}\) yields the displayed formula. Because \(\SeedSplit{\sigma}\) lands in \(\SeedSupportShell{\sigma}\subset\SeedShell{\sigma}\), its lifted image lies in \(\GermSupportShell{\sigma}\) by Proposition \ref{prop:support}. Projecting the resulting point in \(\GermShell{\sigma}\) recovers exactly \(\SeedSplit{\sigma}\), so the projected split identity holds automatically.
\end{proof}

\begin{proposition}[Canonical exact zero-response realization and kernel package]
\label{prop:kernel}
Fix a minimal germ zero-shell transport nucleus \(\GermNucleus\). Then the omitted exact germ zero-response realization and the fixed-data solvability / coercive package are determined canonically by the zero-core diffeomorphism of Proposition \ref{prop:zero}. In particular, one obtains the unique exact germ zero-response realization
\[
\GermEmbed{\sigma}
\equiv
\bigl(\GermProj{\sigma}\big|_{\GermZeroCore{\sigma}}\bigr)^{-1}
\circ
\SeedEmbed{\sigma}
\]
and the unique transported kernel / coercive package on \(\GermZeroCore{\sigma}\) projecting to \(\SeedKernel{\sigma},\SeedCoercive{\sigma}\).
\end{proposition}

\begin{proof}
Because the zero-core projection is a diffeomorphism by Proposition \ref{prop:zero}, the fixed seed exact zero-response realization \(\SeedEmbed{\sigma}\) lifts uniquely to \(\GermZeroCore{\sigma}\) by composition with its inverse. The identity
\[
\SeedEmbed{\sigma}
=
\GermProj{\sigma}\circ\GermEmbed{\sigma}
\]
is then immediate. The same zero-core diffeomorphism transports the fixed seed kernel chart and seed coercive package to a unique fixed-data solvability / coercive package on \(\GermZeroCore{\sigma}\). Any different choice would project to different seed data and would therefore violate Definition \ref{def:core}(4).
\end{proof}

\begin{proposition}[Unique benchmark-faithful core completion]
\label{prop:completion}
Every minimal germ zero-shell transport nucleus \(\GermNucleus\) admits a unique benchmark-faithful core completion in the sense of Definition \ref{def:completion}.
\end{proposition}

\begin{proof}
Definition \ref{def:nucleus} already provides the nucleus data. Proposition \ref{prop:support} canonically supplies the omitted support shell, Proposition \ref{prop:zero} canonically supplies the omitted zero core, Proposition \ref{prop:split} canonically supplies the omitted split and hidden charge, and Proposition \ref{prop:kernel} canonically supplies the omitted exact zero-response realization together with the fixed-data kernel / coercive package. Benchmark faithfulness is already part of the nucleus. Hence a benchmark-faithful core completion exists. Uniqueness follows because each omitted datum is determined by a displayed canonical formula or by transport across a displayed canonical diffeomorphism.
\end{proof}

\begin{proposition}[The nucleus determines the same benchmark-facing consequences]
\label{prop:frontier}
Any two benchmark-faithful current cores with the same minimal germ zero-shell transport nucleus are canonically core-equivalent and induce the same downstream benchmark-facing consequences on the active line.
\end{proposition}

\begin{proof}
By Proposition \ref{prop:completion}, the common nucleus admits only one benchmark-faithful current-core completion. Therefore the two current cores coincide on every omitted current-core constituent after the canonical identifications already built into the completion formulas. In particular, they determine the same germ support shell, the same germ zero core, the same split, the same hidden charge, the same exact germ zero-response realization, and the same fixed-data solvability / coercive package. Since benchmark faithfulness is also part of the nucleus, they preserve the same benchmark outputs \eqref{eq:fixedoutputs}. Thus they induce the same downstream benchmark-facing consequences.
\end{proof}

\begin{proposition}[The benchmark package remains fixed]
\label{prop:benchmark}
Every minimal germ zero-shell transport nucleus preserves the same benchmark one-metric package \eqref{eq:fixedoutputs}. In particular, the operative metric remains exactly \(\CompletedMetric\), the ordinary matter route is unchanged, there is no second operative metric, and the low-energy matter law is not enlarged.
\end{proposition}

\begin{proof}
This is part of benchmark faithfulness in Definition \ref{def:nucleus}(3). The present note removes only internal current-core freedom inside the active derivational line. It does not alter the benchmark exact-closure package.
\end{proof}

\section{Main Theorem}

\begin{theorem}[Observer-relevant germ seed-transport core collapse to a minimal germ zero-shell transport nucleus]
\label{thm:main}
Fix the primitive continuation data of Definition \ref{def:fixed}, the recorded downstream seed package of Definition \ref{def:seed}, and a benchmark-faithful observer-relevant germ seed-transport core in the sense of Definition \ref{def:core}. Then, for each sector \(\sigma\in\{\Car,\Cur,\Hom\}\), the live benchmark-facing burden inside that current core collapses further to the minimal germ zero-shell transport nucleus \(\GermNucleus\) of Definition \ref{def:nucleus}. More precisely:
\begin{enumerate}
    \item every minimal germ zero-shell transport nucleus admits a unique benchmark-faithful current-core completion by Proposition \ref{prop:completion};
    \item the current-core constituents omitted in Definition \ref{def:nucleus} --- namely the germ support shell, germ zero core, germ split, hidden charge, exact germ zero-response realization, and fixed-data solvability / coercive package --- are canonical transport data determined by the nucleus together with the fixed downstream seed package, and therefore do not constitute additional independent benchmark-facing freedom;
    \item any two benchmark-faithful current cores with the same nucleus are canonically core-equivalent and induce the same downstream benchmark-facing consequences by Proposition \ref{prop:frontier};
    \item the benchmark boundary remains fixed by Proposition \ref{prop:benchmark};
    \item consequently the remaining honest burden below the previous germ-core frontier is no longer the full observer-relevant germ seed-transport core as such. What remains is to derive or uniquely force the minimal germ zero-shell transport nucleus itself from still more primitive explicit substrate structure, or else prove a still sharper collapse, sharper necessity result, or sharper uniqueness result strictly inside that nucleus.
\end{enumerate}
\end{theorem}

\begin{proof}
Item (1) is Proposition \ref{prop:completion}. Item (2) is the combined content of Propositions \ref{prop:support}, \ref{prop:zero}, \ref{prop:split}, and \ref{prop:kernel}. Item (3) is Proposition \ref{prop:frontier}. Item (4) is Proposition \ref{prop:benchmark}. Item (5) is the routing consequence of the previous items together with the active safeguard against counting same-output deeper relays as new front-line progress once the live observer-relevant datum is unchanged.
\end{proof}

\begin{corollary}[What no longer counts as the next move]
\label{cor:screened}
After Theorem \ref{thm:main}, none of the following is the default next benchmark-facing move:
\begin{enumerate}
    \item another deeper-origin relay or same-germ relay that leaves the minimal germ zero-shell transport nucleus \(\GermNucleus\) unchanged;
    \item another restatement of the full observer-relevant germ seed-transport core or the full seed-generating substrate germ package that treats the canonically induced omitted data as if they were still independent;
    \item a replay of higher-package consequences that merely preserves the same current nucleus;
    \item a reopening of older screened layers or a manuscript that introduces a second operative metric, enlarges the low-energy matter law, or uses stronger promotion language.
\end{enumerate}
\end{corollary}

\begin{proof}
Items (1)--(3) fail because Theorem \ref{thm:main} shows that the live burden is now smaller than the full current core: once \(\GermNucleus\) is fixed, the omitted current-core constituents are already determined. Item (4) either does not address the live burden at the present frontier or violates benchmark faithfulness.
\end{proof}

\section{Conclusion}

The positive result of this note is narrower than a completed first-principles derivation and sharper than the earlier germ-core collapse. The previous frontier had already shown that the full seed-generating substrate germ package contains benchmark-neutral lower-ladder data. The present theorem then shows that even the current observer-relevant germ seed-transport core still contains benchmark-neutral constituents: the support shell, zero core, split, hidden charge, exact zero-response realization, and fixed-data solvability / coercive package are canonically induced once the minimal germ zero-shell transport nucleus is fixed.

The live burden is therefore cleaner again. At current scope the honest benchmark-facing datum is no longer the full observer-relevant germ seed-transport core \(\GermTransportCore\) but the smaller minimal germ zero-shell transport nucleus \(\GermNucleus\). That is the datum that now requires deeper justification, unless a still smaller theorem-level collapse remains available strictly inside it.

The scope remains disciplined. The benchmark package is unchanged: the one operative metric is still exactly \(\CompletedMetric\), the ordinary matter route is unchanged, there is no second operative metric, and the low-energy matter law is not enlarged. This remains a theorem inside the active derivational program of \TheoryName{}, not a basis for stronger promotion language. The next honest burden is to derive or uniquely force \(\GermNucleus\) from still more primitive explicit substrate structure, or else prove a still sharper collapse, sharper necessity result, or sharper uniqueness result strictly inside it while preserving the same benchmark one-metric package and claim boundary.

\end{document}
